\section{28 сентября 1914 г.}

\lp{20260604_193515915}

Сегодня едем с Лёлей (и конечно берем с собой Луизу Федоровну) навещать
пленных раненных, находящихся в военном госпитале и везем им гостинцев. Л.Ф
будет узнавать про своего брата. Все что бы я ни делала, о чем бы ни думала,
все время я с вами. Я так глубоко чувствую что вы должны сейчас переживать. То,
что я прочла в вашем письме из Мюнхена, все ваши воспоминания о нашей там
жизни, так глубоко потрясло меня и так тронуло, и так захотелось мне быть
сейчас около вас! Что до вопроса о том, не испортили ли вы мне здоровья, жизни
и прочего, то право же, это все в самом благополучном виде, и я чувствую себя
очень нормально и никакие вопросы не встают передо мной. Ничего, ничего вы мне
не испортили. Я говорю вам сущую правду, поверьте мне! И вообще я ничего от вас
не скрываю. Что касается Коли, то он действительно сейчас очень плохо себя
чувствует, не может сочинять и вообще уж ничего еще не может. Даже не читает
ничего кроме газет. И он тоже все о вас думает и страдает вместе с вами. Мы,
конечно, не верим всему, что читаем, но и вы не должны верить всему. Никакого
сомнения быть не может, что н. тоже не пишут про все, что они делают.

Почему вы пишите, что война не помешает вам осуществить союз с Ядвигой? Разве
сейчас не настоящий момент ей показать мое отношение к вам? Ведь теперь-то ей и
ее маменьке ясно, что вы ничего не можете поделать. Она бы теперь
просто-напросто должна была поселиться с вами в Швейцарии. Или же она совсем
только и думает о своем устройстве. Единственное, что сейчас по-моему, может
явиться оправданием -- это если вам сейчас

\lp{20260604_193529896}

\noindent
надо и полезно быть одному. Но если Ядвигины мотивы -- ожидание формального
развода, то право же, она очень среднее существо. Ну да Бог с ней!

Письмо Смирнову я отнесла. Фридрихам послала. Жиляев призван, но еще не на
войне, а где-то тут, под Москвой заведует вещевым складом. Степун призван и
находится сейчас в Сибири. Кожебаткин уже приперт к стене (по словам
Киселева). Н. П. говорит, что как раз когда вы писали мне об этом, он очевидно
на расстоянии почувствовал ваше это желание и деятельно принялся за
Кожебаткина. На днях он должен что-то вернуть, а то подвергнется уголовному
взысканию. Сабурова переводят в Москву, а пока что он в Калуге. Он так зорошо
обучил свою дружину, что его не отправили с ней на войну, а оставили, чтобы
обучать еще дружины.

Совершенно то же самое, что и вы говорите о братоубийстве, говорит и Коля. Ему
тоже так тяжело, что право, невозможно смотреть на него. Всех, всех людей так
жалко! Дай Бог скорей бы все это кончилось! Все вас обнимают и очень любят.
Постоянно все говорим о вас. Пишу каждый день понемногу, но отсылаю раза два,
три в неделю. Числа должны получаться непрерывные. Обнимаю вас, мой дорогой с
бесконечной любовью! Ваша Анюта.
